% Generated by publication/build_sections.py; do not edit by hand.
\paragraph{Informal verdict.}
\textbf{C --- Substantial progress}, with \textbf{high confidence in the central classification}.

\paragraph{Lean coverage (partial).}
Lean proves the structural periodic-point and asymptotic classification, the power-of-two criterion, the complete binary 15-cycle for n=5, and the explicit eigenray tending to zero. The Section 6 example and the general basin-of-attraction problem are not formalized or claimed solved.

\subsection{Audited informal answer}
\subsection{Partial result}

I give below:

\begin{enumerate}
\item a complete classification of all possible periodic attractors for every \(n\);
\item a proof that every real orbit is asymptotic, phase by phase, to one of them;
\item an exact finite algorithm for rational data;
\item a complete description of the possible asymptotic cycles for \(n=5\);
\item an explicit nonconstant \(5\)-tuple converging to \(0\);
\item an exact algebraic \(5\)-tuple which approaches the nonzero \(15\)-cycle but never enters any periodic orbit.
\end{enumerate}

What remains unresolved by this analysis is a direct criterion, for an arbitrary exact real \(X\), deciding which attractor basin contains \(X\) and computing the limiting scale without iterating indefinitely.

\subsubsection{1. Faithful restatement}

Let \(n\ge 3\). For

\[
X=(x_1,\ldots,x_n)\in\mathbb R_+^n,
\]

define

\[
P(X)=
\bigl(
|x_1-x_2|,
|x_2-x_3|,
\ldots,
|x_{n-1}-x_n|,
|x_n-x_1|
\bigr).
\]

Study the orbit

\[
X,\ P(X),\ P^2(X),\ldots
\]

for a prescribed starting vector \(X\).

The editorial note says that \(n=4\) is understood: every orbit eventually reaches \(0\), except for one algebraic family, up to elementary transformations, represented after one iteration by

\[
(1,a,a^2,a^3),
\]

where \(a-1\) is the real root of

\[
u^3+2u^2-2=0.
\]

Equivalently, \(a\) satisfies

\[
a^3-a^2-a-1=0.
\]

The note explicitly welcomes any partial result encompassing \(n=5\). 

There are two minor notational points.

\begin{itemize}
\item Here \(\mathbb R_+\) must mean \([0,\infty\)), since \(P(X)\) can have zero coordinates.
\item The phrase ``up to a simple transformation'' in the editorial note is informal and does not specify the transformations, but this does not affect the formal problem.
\end{itemize}

Throughout, write

\[
X^{(q)}=P^q(X),\qquad
M_q=|X^{(q)}|_\infty.
\]

\medskip\hrule\medskip

\subsection{2. Main classification theorem}

Write

\[
n=hm,\qquad h=2^{\nu_2(n)},\qquad m\ \text{odd}.
\]

Thus \(h\) is the largest power of \(2\) dividing \(n\).

For a binary vector
\(\varepsilon=(\varepsilon_1,\ldots,\varepsilon_n)\in{0,1}^n\),
put

\[
E_\varepsilon(z)=
\sum_{j=1}^{n}\varepsilon_jz^{j-1}
\in\mathbb F_2[z].
\]

Define

\[
\mathcal S_n=
\left\{
\varepsilon\in{0,1}^n:
1+z^h\mid E_\varepsilon(z)
\text{ in }\mathbb F_2[z]
\right\}.
\]

Equivalently,

\[
\boxed{
\varepsilon\in\mathcal S_n
\iff
\sum_{k=0}^{m-1}\varepsilon_{r+kh}\equiv0\pmod2
\quad
(1\le r\le h).
}
\tag{2.1}
\]

Thus \(\mathcal S_n\) is given by \(h\) independent parity conditions and

\[
|\mathcal S_n|=2^{n-h}.
\]

\subsubsection{Theorem 2.1}

For every \(X\in\mathbb R_+^n\):

\begin{enumerate}
\item The sequence \(M_q\) is nonincreasing, and therefore
\end{enumerate}

\[
   c(X):=\lim_{q\to\infty}M_q
\]

   exists.

\begin{enumerate}
\item There exists a periodic binary sequence
\end{enumerate}
   \((\varepsilon^{(q)})_{q\ge0}\), with every
   \(\varepsilon^{(q)}\in\mathcal S_n\) and

\[
   \varepsilon^{(q+1)}=P(\varepsilon^{(q)}),
\]

   such that

\[
   \boxed{
   \lim_{q\to\infty}
   \left\lvert 
   P^q(X)-c(X)\varepsilon^{(q)}
   \right\rvert _\infty=0.
   }
   \tag{2.2}
\]

\begin{enumerate}
\item The periodic points of \(P\) in \(\mathbb R_+^n\) are exactly
\end{enumerate}

\[
   \boxed{
   c\varepsilon,
   \qquad c\ge0,\quad \varepsilon\in\mathcal S_n.
   }
   \tag{2.3}
\]

\begin{enumerate}
\item The map \(P\) permutes the finite set \(\mathcal S_n\). Hence all possible periodic cycles can be computed by taking the cycle decomposition of this permutation.
\end{enumerate}

\begin{enumerate}
\item Every real orbit converges to \(0\) if and only if \(n\) is a power of \(2\).
\end{enumerate}

An equivalent limit-point and periodic-orbit classification appears in the classical work of Misiurewicz and Schinzel; the proof below is self-contained and also makes the phase-by-phase convergence explicit. ([hrj.episciences.org][1])

\medskip\hrule\medskip

\subsection{3. Proof of Theorem 2.1}

\subsubsection{3.1. Monotonicity}

If \(Y=(y_1,\ldots,y_n)\in[0,M]^n\), then

\[
|y_i-y_{i+1}|\le M.
\]

Consequently,

\[
|P(Y)|*\infty\le |Y|*\infty.
\]

Thus \(M_{q+1}\le M_q\), proving the existence of \(c(X)\).

We shall need the following strict form.

\subsubsection{3.2. Strict decrease unless all entries are extremal}

For \(u_0,\ldots,u_k\in[0,M]\), define recursively

\[
\Delta_0(u_0)=u_0
\]

and

\[
\Delta_k(u_0,\ldots,u_k)
=
%
\left\lvert 
\Delta_{k-1}(u_0,\ldots,u_{k-1})
--------------------------------
%
\Delta_{k-1}(u_1,\ldots,u_k)
\right\rvert .
\]

Then, with cyclic indices,

\[
\bigl(P^k(Y)\bigr)_i
=
%
\Delta_k(y_i,y_{i+1},\ldots,y_{i+k}).
\tag{3.1}
\]

\paragraph{Lemma 3.1}

Suppose \(u_0,\ldots,u_k\in[0,M]\). If

\[
\Delta_k(u_0,\ldots,u_k)=M,
\]

then

\[
u_j\in{0,M}
\qquad(0\le j\le k).
\]

\paragraph{Proof}

First observe the following endpoint fact.

If

\[
u_0,\ldots,u_{k-1}\in{0,M},
\qquad u_k=t\in[0,M],
\]

then

\[
\Delta_k(u_0,\ldots,u_k)\in{t,M-t}.
\tag{3.2}
\]

Indeed, after taking one row of absolute differences, all entries except possibly the last are again in (\{0,M\}), while the last is either \(t\) or \(M-t\). Induction proves (3.2). The analogous assertion holds when the exceptional entry is the first one.

We now prove the lemma by induction on \(k\). The case \(k=0\) is immediate.

Set

\[
A=\Delta_{k-1}(u_0,\ldots,u_{k-1}),\qquad
B=\Delta_{k-1}(u_1,\ldots,u_k).
\]

Both \(A\) and \(B\) belong to ([0,M]). Since

\[
|A-B|=M,
\]

we must have

\[
{A,B}={0,M}.
\]

Suppose, for example, that \(A=M\). By the induction hypothesis,

\[
u_0,\ldots,u_{k-1}\in{0,M}.
\]

Since \(B=0\), the endpoint observation applied to
\(u_1,\ldots,u_k\) implies that \(u_k\in{0,M}\).
The other case is symmetric. \(\square\)

\paragraph{Corollary 3.2}

If \(Y\in[0,M]^n\) has at least one coordinate strictly between \(0\) and \(M\), then

\[
\boxed{
|P^{n-1}(Y)|_\infty<M.
}
\tag{3.3}
\]

\paragraph{Proof}

Every coordinate of \(P^{n-1}(Y)\) is an (\(n-1\))-st absolute difference involving all \(n\) original coordinates once. Lemma 3.1 shows that none of these coordinates can equal \(M\). \(\square\)

\medskip\hrule\medskip

\subsubsection{3.3. Every limit point is binary up to scale}

Let \(\Omega(X)\) be the set of limit points of the orbit.

Take \(Y\in\Omega(X)\). There is a sequence \(q_j\to\infty\) such that

\[
X^{(q_j)}\longrightarrow Y.
\]

Since \(M_{q_j}\to c(X)\),

\[
|Y|_\infty=c(X).
\]

Suppose \(c(X)>0\) and that some coordinate of \(Y\) lies strictly between
\(0\) and \(c(X)\). Corollary 3.2 gives

\[
|P^{n-1}(Y)|_\infty<c(X).
\]

On the other hand, continuity of \(P\) yields

\[
P^{n-1}(X^{(q_j)})
\longrightarrow
P^{n-1}(Y),
\]

while

\[
\left\lvert P^{n-1}(X^{(q_j)})\right\rVert_\infty
=M_{q_j+n-1}\longrightarrow c(X),
\]

a contradiction.

Therefore every limit point has the form

\[
Y=c(X)\varepsilon,
\qquad
\varepsilon\in{0,1}^n.
\tag{3.4}
\]

When \(c(X)=0\), the orbit itself converges to \(0\).

\medskip\hrule\medskip

\subsubsection{3.4. Algebra over \(\mathbb F_2\)}

On binary vectors, absolute difference agrees with addition modulo \(2\):

\[
|a-b|=a+b\pmod2
\qquad(a,b\in{0,1}).
\]

Let \(L:\mathbb F_2^n\to\mathbb F_2^n\) denote the resulting linear map

\[
(L\varepsilon)_i=\varepsilon_i+\varepsilon_{i+1}.
\]

Identify \(\mathbb F_2^n\) with

\[
R_n=\mathbb F_2[z]/(z^n-1)
\]

by \(\varepsilon\mapsto E_\varepsilon(z)\). Then

\[
E_{L\varepsilon}(z)
=
%
z^{-1}(1+z)E_\varepsilon(z)
\quad\text{in }R_n.
\tag{3.5}
\]

Since multiplication by \(z\) is invertible in \(R_n\), the image of \(L^h\) is the ideal generated by

\[
(1+z)^h=1+z^h,
\]

where we used that \(h\) is a power of \(2\).

Moreover \(1+z^h\mid z^n-1\) in \(\mathbb F_2[z]\). It follows that

\[
\operatorname{im}L^h=\mathcal S_n.
\tag{3.6}
\]

Reducing \(E_\varepsilon\) modulo \(1+z^h\) gives

\[
E_\varepsilon(z)
\equiv
\sum_{r=1}^{h}
\left(
\sum_{k=0}^{m-1}\varepsilon_{r+kh}
\right)z^{r-1},
\]

which proves the parity formulation (2.1).

\medskip\hrule\medskip

\subsubsection{3.5. Why limit points belong to \(\mathcal S_n\)}

The omega-limit set satisfies

\[
P^k(\Omega(X))=\Omega(X)
\qquad(k\ge1).
\tag{3.7}
\]

Indeed, the inclusion from left to right follows by continuity. Conversely, if

\[
Y=\lim_j X^{(q_j)},
\]

take a convergent subsequence of \(X^{(q_j-k)}\); its limit \(Z\in\Omega(X)\) satisfies \(P^k(Z)=Y\).

Now let \(Y=c\varepsilon\in\Omega(X)\), with \(c>0\). By (3.7), there exists \(c\delta\in\Omega(X)\) such that

\[
P^h(c\delta)=c\varepsilon.
\]

Hence

\[
\varepsilon=L^h\delta\in\operatorname{im}L^h=\mathcal S_n.
\]

Thus

\[
\Omega(X)\subseteq c(X)\mathcal S_n.
\tag{3.8}
\]

\medskip\hrule\medskip

\subsubsection{3.6. \(P\) is a permutation of \(\mathcal S_n\)}

Since

\[
\mathcal S_n=\operatorname{im}L^h,
\]

we have \(L(\mathcal S_n)\subseteq\mathcal S_n\).

It remains to prove injectivity. The kernel of \(L\) consists of the two constant binary vectors:

\[
\ker L={0,\mathbf 1},
\qquad
\mathbf 1=(1,\ldots,1).
\]

We claim that \(\mathbf 1\notin\mathcal S_n\). Its polynomial is

\[
J_n(z)=1+z+\cdots+z^{n-1}
=\frac{z^n+1}{z+1}
\quad\text{in characteristic }2.
\]

Because \(n=hm\) with \(m\) odd,

\[
z^n+1=(z^m+1)^h.
\]

The root \(z=1\) has multiplicity exactly one in \(z^m+1\), since

\[
\frac{d}{dz}(z^m+1)\Big|_{z=1}
=m\equiv1\pmod2.
\]

Thus \(z=1\) has multiplicity \(h\) in \(z^n+1\), and multiplicity \(h-1\) in \(J_n\). Therefore

\[
(1+z)^h=1+z^h\nmid J_n(z).
\]

So \(\mathbf 1\notin\mathcal S_n\).

If \(L\varepsilon=L\delta\) with
\(\varepsilon,\delta\in\mathcal S_n\), then

\[
\varepsilon+\delta\in\ker L.
\]

It cannot equal \(\mathbf1\), because \(\mathcal S_n\) is a vector space and \(\mathbf1\notin\mathcal S_n\). Hence \(\varepsilon=\delta\).

Thus \(L\), and therefore \(P\), permutes \(\mathcal S_n\).

It also follows that the binary periodic points are precisely the elements of \(\mathcal S_n\). Indeed:

\begin{itemize}
\item every element of \(\mathcal S_n\) is periodic because \(L\) permutes this finite set;
\item if \(L^r\varepsilon=\varepsilon\), choose a multiple \(s\ge h\) of \(r\); then
\end{itemize}

\[
  \varepsilon=L^s\varepsilon
  =L^h(L^{s-h}\varepsilon),
\]

  so \(\varepsilon\in\operatorname{im}L^h=\mathcal S_n\).

\medskip\hrule\medskip

\subsubsection{3.7. The orbit shadows one periodic cycle}

We know that all limit points lie in the finite set

\[
F=c(X)\mathcal S_n.
\]

Consequently,

\[
\operatorname{dist}(X^{(q)},F)\longrightarrow0.
\tag{3.9}
\]

Otherwise a subsequence staying a positive distance from \(F\) would have a limit point outside \(F\).

Assume \(c(X)>0\). Distinct points of \(F\) are separated in the supremum norm by at least \(c(X)\).

The map \(P\) is \(2\)-Lipschitz:

\[
|P(U)-P(V)|*\infty
\le2|U-V|*\infty.
\tag{3.10}
\]

For all sufficiently large \(q\), there is therefore a unique \(f_q\in F\) with

\[
|X^{(q)}-f_q|_\infty<\frac{c(X)}{10}.
\]

Furthermore,

\[
\begin{aligned}
|P(f_q)-f_{q+1}|*\infty
&\le
|P(f_q)-P(X^{(q)})|*\infty
+
|X^{(q+1)}-f_{q+1}|_\infty\\
&<\frac{3c(X)}{10}.
\end{aligned}
\]

Since two distinct elements of \(F\) are at distance at least \(c(X)\), this forces

\[
f_{q+1}=P(f_q).
\]

Hence, from some point onward, the nearest elements \(f_q\) follow one periodic cycle. Together with (3.9), this proves (2.2).

\medskip\hrule\medskip

\subsubsection{3.8. All real periodic points}

Let \(Y\) be periodic. The sequence \(|P^q(Y)|_\infty\) is both periodic and nonincreasing, hence constant.

If some coordinate of \(Y\) were strictly between \(0\) and
\(M=|Y|_\infty\), Corollary 3.2 would give

\[
|P^{n-1}(Y)|_\infty<M,
\]

a contradiction. Thus

\[
Y=M\varepsilon
\]

for a binary vector \(\varepsilon\). The binary classification proves
\(\varepsilon\in\mathcal S_n\).

Conversely, \(M\varepsilon\) is periodic whenever
\(\varepsilon\in\mathcal S_n\), by positive homogeneity:

\[
P(M\varepsilon)=M P(\varepsilon).
\]

This proves (2.3).

Finally, if \(n\) is a power of \(2\), then \(h=n\), so the only degree-\(<n\) polynomial divisible by \(1+z^n\) is zero. Hence

\[
\mathcal S_n={0}
\]

and every orbit converges to \(0\).

If \(n\) is not a power of \(2\), then \(h<n\), so

\[
|\mathcal S_n|=2^{n-h}>1,
\]

and nonzero periodic orbits exist. \(\square\)

\medskip\hrule\medskip

\subsection{4. Complete asymptotic description for \(n=5\)}

For \(n=5\),

\[
h=1,\qquad m=5.
\]

Condition (2.1) becomes

\[
\varepsilon_1+\cdots+\varepsilon_5\equiv0\pmod2.
\]

Thus

\[
\boxed{
\mathcal S_5=
{\varepsilon\in{0,1}^5:
\varepsilon\text{ has even Hamming weight}}.
}
\tag{4.1}
\]

There are \(16\) such vectors.

\subsubsection{Proposition 4.1}

The zero vector is fixed, and the other \(15\) vectors of \(\mathcal S_5\) form one cycle of exact length \(15\).

\paragraph{Proof}

Set

\[
H(z)=1+z+z^2+z^3+z^4.
\]

This polynomial is irreducible over \(\mathbb F_2\):

\begin{itemize}
\item it has no root in \(\mathbb F_2\);
\item if it factored without a linear factor, it would be a product of two irreducible quadratics;
\item the only monic irreducible quadratic over \(\mathbb F_2\) is
\end{itemize}
  \(z^2+z+1\), whose square is \(z^4+z^2+1\ne H(z)\).

Consequently,

\[
K=\mathbb F_2[z]/(H)
\]

is the field with \(16\) elements.

The space \(\mathcal S_5=(1+z)R_5\) is naturally isomorphic, as a vector space, to \(K\). Under this identification, the binary Ducci map acts as multiplication by

\[
\mu=1+\alpha,
\]

where \(\alpha\) is a nontrivial fifth root of unity in \(K\).

Since

\[
1+\alpha+\alpha^2+\alpha^3+\alpha^4=0,
\]

we get

\[
\mu^3=(1+\alpha)^3
=1+\alpha+\alpha^2+\alpha^3
=\alpha^4.
\]

Thus \(\mu^3\) has order \(5\). Since \(K^\times\) has order \(15\), the order of \(\mu\) is either \(5\) or \(15\).

But

\[
\mu^5=(1+\alpha)^5
=(1+\alpha^4)(1+\alpha)
=\alpha+\alpha^4.
\]

If \(\mu^5=1\), then

\[
1+\alpha+\alpha^4=0.
\]

Adding \(1+\alpha+\alpha^2+\alpha^3+\alpha^4=0\) gives

\[
\alpha^2+\alpha^3=\alpha^2(1+\alpha)=0,
\]

which is impossible. Therefore \(\mu\) has order \(15\).

Multiplication by \(\mu\) consequently acts transitively on the \(15\) nonzero elements of \(K\). \(\square\)

One explicit orientation of the cycle is

\[
\begin{aligned}
e_0={}&(1,1,0,0,0),\\
e_1={}&(0,1,0,0,1),\\
e_2={}&(1,1,0,1,1),\\
e_3={}&(0,1,1,0,0),\\
e_4={}&(1,0,1,0,0),\\
e_5={}&(1,1,1,0,1),\\
e_6={}&(0,0,1,1,0),\\
e_7={}&(0,1,0,1,0),\\
e_8={}&(1,1,1,1,0),\\
e_9={}&(0,0,0,1,1),\\
e_{10}={}&(0,0,1,0,1),\\
e_{11}={}&(0,1,1,1,1),\\
e_{12}={}&(1,0,0,0,1),\\
e_{13}={}&(1,0,0,1,0),\\
e_{14}={}&(1,0,1,1,1),
\end{aligned}
\]

with

\[
P(e_q)=e_{q+1},\qquad e_{15}=e_0.
\]

\subsubsection{Corollary 4.2}

For every \(X\in\mathbb R_+^5\), exactly one of the following holds.

\paragraph{Zero-attractor case}

\[
\boxed{P^q(X)\longrightarrow0.}
\]

\paragraph{Nonzero-attractor case}

There are \(c>0\) and a phase \(r\in{0,\ldots,14}\) such that

\[
\boxed{
\left\lvert 
P^q(X)-c,e_{q+r\pmod{15}}
\right\rvert _\infty
\longrightarrow0.
}
\tag{4.2}
\]

Thus the possible asymptotic shapes for \(n=5\) are completely classified.

\medskip\hrule\medskip

\subsection{5. A nonconstant \(5\)-tuple converging to zero}

The zero-attractor case is not confined to constant vectors.

Let \(\rho\in(0,1)\) be a root of

\[
\rho^4+3\rho^3+2\rho^2-2\rho-3=0.
\tag{5.1}
\]

Such a root exists because the polynomial is negative at \(0\) and positive at \(1\). Numerically,

\[
\rho\approx0.9275619755.
\]

Define

\[
V=
\left(
1,
1-\rho,
1-\rho^2,
(1+\rho)(1-\rho^2),
(1+\rho)^2(1-\rho^2)
\right).
\]

The identity

\[
(1+\rho)^3(1-\rho^2)-1
=
%
-\rho\bigl(\rho^4+3\rho^3+2\rho^2-2\rho-3\bigr)
\]

shows that

\[
(1+\rho)^3(1-\rho^2)=1.
\]

The coordinates satisfy

\[
0<v_2<v_3<v_4<v_5<v_1.
\]

Moreover,

\[
\begin{aligned}
v_1-v_2&=\rho v_1,\\
v_3-v_2&=\rho v_2,\\
v_4-v_3&=\rho v_3,\\
v_5-v_4&=\rho v_4,\\
v_1-v_5&=\rho v_5.
\end{aligned}
\]

Therefore

\[
\boxed{P(V)=\rho V.}
\]

Consequently,

\[
P^q(V)=\rho^qV\longrightarrow0,
\]

but \(P^q(V)\ne0\) for every finite \(q\).

\medskip\hrule\medskip

\subsection{6. An exact non-eventually-periodic \(n=5\) orbit with nonzero attractor}

This is the more delicate phenomenon. Earlier work explicitly asked whether the \(n=5\) situation permits an orbit which approaches a nontrivial periodic cycle without ever becoming periodic. ([fq.math.ca][2])

The following gives such an orbit exactly.

\subsubsection{Theorem 6.1}

There exists \(X\in\mathbb R_+^5\) such that

\[
\left\lvert P^q(X)-e_{q\bmod15}\right\rvert _\infty\longrightarrow0,
\]

but no \(P^q(X)\) is periodic.

\paragraph{Construction}

Let

\[
p(t)=t^3-19t^2+139t-1.
\]

We have

\[
p\left(\frac1{139}\right)
=
%
-\frac{2640}{2685619}<0,
\]

and

\[
p\left(\frac1{138}\right)
=
%
\frac{16423}{2628072}>0.
\]

Also \(p'(t)>0\) on this interval. Let

\[
\lambda\in\left(\frac1{139},\frac1{138}\right)
\]

be its unique root. Numerically,

\[
\lambda\approx0.007201330580038611.
\]

Define

\[
w_0=
\begin{pmatrix}
19960\\
83\lambda^2-424\lambda+18633\\
22\lambda^2-1916\lambda+6502\\
112\lambda^2-2496\lambda+9512\\
7\lambda^2-156\lambda+4337
\end{pmatrix}.
\tag{6.1}
\]

For a sign word
\(\sigma=(\sigma_1,\ldots,\sigma_5)\in{+,-}^5\), let \(B_\sigma\) be the matrix defined by

\[
(B_\sigma v)_i=s_i(v_i-v_{i+1}),
\qquad
s_i=
\begin{cases}
1,&\sigma_i=+,\\
-1,&\sigma_i=-,
\end{cases}
\]

with cyclic indices.

Use the following \(15\) sign words:

\[
\begin{array}{c|c@{\qquad}c|c@{\qquad}c|c}
q&\sigma_q&q&\sigma_q&q&\sigma_q\\ \hline
0&++-+-&5&+-+--&10&+-+-+\\
1&-+--+&6&--++-&11&-+-++\\
2&++--+&7&-+-++&12&+++--\\
3&--+++&8&++-+-&13&+--+-\\
4&+-+--&9&+---+&14&+--+-
\end{array}
\tag{6.2}
\]

Define recursively

\[
w_{q+1}=B_{\sigma_q}w_q
\qquad(0\le q<15).
\]

The product is

\[
A=B_{\sigma_{14}}\cdots B_{\sigma_0}
=
%
\begin{pmatrix}
2&-1&-15& 3&11\\
-9&10&-13& 5& 7\\
11&-11&10&-7&-3\\
1&-1& 9&-4&-5\\
1&-1&-3& 1& 2
\end{pmatrix}.
\tag{6.3}
\]

If \(w_0(t)\) denotes the vector in (6.1) with \(\lambda\) replaced by \(t\), direct multiplication gives

\[
A w_0(t)-t w_0(t)
=
%
p(t)
\begin{pmatrix}
0\-83\-22\-112\-7
\end{pmatrix}.
\tag{6.4}
\]

Hence

\[
w_{15}=Aw_0=\lambda w_0.
\tag{6.5}
\]

The finite exact verification following the proof establishes two further facts:

\[
w_{q,i}>0
\qquad(0\le q<15,\ 1\le i\le5),
\tag{6.6}
\]

and whenever \(e_{q,i}=e_{q,i+1}\), the corresponding sign in \(\sigma_q\) is the sign of

\[
w_{q,i}-w_{q,i+1}.
\tag{6.7}
\]

Let

\[
K=
\max_{\substack{0\le q<15\\1\le i\le5}}
|w_{q,i}-w_{q,i+1}|,
\qquad
W=
\max_{\substack{0\le q<15\\1\le i\le5}}
w_{q,i}.
\]

Choose \(\varepsilon>0\) so small that

\[
\varepsilon K<\frac12,
\qquad
\varepsilon W<\frac12.
\]

Set

\[
X=e_0+\varepsilon w_0.
\tag{6.8}
\]

\paragraph{Proof of the claimed dynamics}

For every \(q\), consider

\[
e_q+\delta w_q
\]

with \(0<\delta\le\varepsilon\).

If \(e_{q,i}\ne e_{q,i+1}\), their difference is \(\pm1\), and
\(\delta K<1/2\) ensures that the perturbation cannot change its sign.

If \(e_{q,i}=e_{q,i+1}\), condition (6.7) determines the sign of the perturbed difference.

Therefore

\[
P(e_q+\delta w_q)
=
%
e_{q+1}+\delta w_{q+1}.
\tag{6.9}
\]

Using \(w_{15}=\lambda w_0\) and induction gives the exact formula

\[
\boxed{
P^{15k+q}(X)
=
%
e_q+\varepsilon\lambda^k w_q,
\qquad
k\ge0,\quad0\le q<15.
}
\tag{6.10}
\]

Since \(0<\lambda<1\),

\[
P^{15k+q}(X)\longrightarrow e_q.
\]

Thus the orbit approaches the nonzero \(15\)-cycle.

On the other hand, every coordinate in (6.10) is strictly positive by (6.6). Every nonzero periodic point for \(n=5\) is a scalar multiple of an even-weight binary vector, and hence has at least one zero coordinate. Thus no iterate in (6.10) is periodic.

This proves Theorem 6.1. \(\square\)

\subsubsection{Exact finite verifier}

This is not a floating-point experiment. All inequalities are certified using rational interval arithmetic on

\[
\lambda\in\left[\frac1{139},\frac1{138}\right]
\]

and reductions modulo \(p(\lambda)=0\).

\begin{verbatim}
from fractions import Fraction
import sympy as sp

t = sp.symbols("t")
p = t**3 - 19*t**2 + 139*t - 1
lo, hi = Fraction(1, 139), Fraction(1, 138)

def peval(x):
    return x**3 - 19*x**2 + 139*x - 1

# Existence and uniqueness of the small positive root.
assert peval(lo) < 0 < peval(hi)
assert 3*lo*lo - 38*hi + 139 > 0

cycle = [
    (1,1,0,0,0),
    (0,1,0,0,1),
    (1,1,0,1,1),
    (0,1,1,0,0),
    (1,0,1,0,0),
    (1,1,1,0,1),
    (0,0,1,1,0),
    (0,1,0,1,0),
    (1,1,1,1,0),
    (0,0,0,1,1),
    (0,0,1,0,1),
    (0,1,1,1,1),
    (1,0,0,0,1),
    (1,0,0,1,0),
    (1,0,1,1,1),
]

words = [
    "++-+-", "-+--+", "++--+", "--+++",
    "+-+--", "+-+--", "--++-", "-+-++",
    "++-+-", "+---+", "+-+-+", "-+-++",
    "+++--", "+--+-", "+--+-",
]

def branch_matrix(word):
    B = sp.zeros(5)
    for i, symbol in enumerate(word):
        s = 1 if symbol == "+" else -1
        B[i, i] = s
        B[i, (i + 1) % 5] = -s
    return B

def reduce_mod_p(expr):
    return sp.rem(
        sp.Poly(sp.expand(expr), t),
        sp.Poly(p, t),
    ).as_expr()

def reduce_vector(v):
    return sp.Matrix([
        sp.expand(reduce_mod_p(x)) for x in v
    ])

w0 = sp.Matrix([
    19960,
    83*t**2 - 424*t + 18633,
    22*t**2 - 1916*t + 6502,
    112*t**2 - 2496*t + 9512,
    7*t**2 - 156*t + 4337,
])

ws = [w0]
A = sp.eye(5)

for q, word in enumerate(words):
    # Check the binary cycle and fixed signs on unequal edges.
    assert tuple(
        abs(cycle[q][i] - cycle[q][(i + 1) % 5])
        for i in range(5)
    ) == cycle[(q + 1) % 15]

    for i, symbol in enumerate(word):
        if cycle[q][i] != cycle[q][(i + 1) % 5]:
            s = 1 if symbol == "+" else -1
            assert s * (
                cycle[q][i] - cycle[q][(i + 1) % 5]
            ) == 1

    B = branch_matrix(word)
    A = B * A
    ws.append(reduce_vector(B * ws[-1]))

A_expected = sp.Matrix([
    [ 2, -1, -15,  3, 11],
    [-9, 10, -13,  5,  7],
    [11,-11,  10, -7, -3],
    [ 1, -1,   9, -4, -5],
    [ 1, -1,  -3,  1,  2],
])

assert A == A_expected
assert all(reduce_mod_p(x) == 0 for x in A*w0 - t*w0)
assert all(reduce_mod_p(x) == 0 for x in ws[15] - t*w0)

def coefficients(expr):
    poly = sp.Poly(expr, t, domain=sp.QQ)
    out = []
    for degree in range(3):
        c = poly.coeff_monomial(t**degree)
        out.append(Fraction(int(c.p), int(c.q)))
    return out

def interval(expr):
    powers = [
        (Fraction(1), Fraction(1)),
        (lo, hi),
        (lo*lo, hi*hi),
    ]

    lower = Fraction(0)
    upper = Fraction(0)

    for coefficient, (a, b) in zip(coefficients(expr), powers):
        if coefficient >= 0:
            lower += coefficient * a
            upper += coefficient * b
        else:
            lower += coefficient * b
            upper += coefficient * a

    return lower, upper

lower_bounds = []

for q in range(15):
    for i in range(5):
        # Positivity of every perturbation coordinate.
        lower_bounds.append(interval(ws[q][i])[0])

        # Correct sign whenever the binary edge is zero.
        j = (i + 1) % 5
        if cycle[q][i] == cycle[q][j]:
            s = 1 if words[q][i] == "+" else -1
            lower_bounds.append(
                interval(s * (ws[q][i] - ws[q][j]))[0]
            )

minimum = min(lower_bounds)

print("Return matrix:")
print(A)
print("Certified minimum:", minimum, float(minimum))

assert minimum == Fraction(28286630, 1333149)
assert minimum > 21
\end{verbatim}

The certified lower bound is

\[
\frac{28286630}{1333149}>21,
\]

so all required inequalities are strict with a wide margin.

\medskip\hrule\medskip

\subsection{7. Exact algorithm for rational starting vectors}

Suppose

\[
X=(x_1,\ldots,x_n)\in\mathbb Q_+^n.
\]

Choose a common denominator \(D\) and put

\[
A=DX\in\mathbb Z_+^n.
\]

Positive homogeneity gives

\[
P^q(X)=\frac1D P^q(A).
\]

Let

\[
R=\max_i A_i-\min_i A_i.
\]

Every coordinate of \(P(A)\) lies in \({0,\ldots,R}\), and the same is true for every later iterate. Thus after the first step there are at most

\[
(R+1)^n
\]

possible states.

The following exact algorithm therefore terminates.

\begin{verbatim}
INPUT: rational vector X of length n

1. Let D be the least common multiple of all denominators.
2. Set A := D X, an integer vector.
3. Set Y := P(A).
4. Initialize an empty dictionary seen.
5. For q = 1,2,3,...:
       if Y is already in seen:
           μ := seen[Y]
           ℓ := q - μ
           return transient length μ,
                  cycle length ℓ,
                  exact cycle divided by D
       seen[Y] := q
       Y := P(Y)
\end{verbatim}

Each iteration takes \(O(n)\) integer operations on integers having
\(O(\log R)\) bits.

The crude worst-case bounds are therefore

\[
O\bigl(n(R+1)^n\bigr)
\]

arithmetic operations and

\[
O\bigl(n(R+1)^n\log R\bigr)
\]

bit storage. Floyd's cycle-finding method reduces the memory to \(O(n\log R)\) bits, at the cost of rerunning part of the orbit to output the cycle.

The same procedure applies whenever there exist \(a\ne0\) and \(b\in\mathbb R\) such that

\[
a(X+b\mathbf1)\in\mathbb Q^n,
\]

because

\[
P(X+b\mathbf1)=P(X),\qquad
P(aX)=|a|P(X).
\]

\subsubsection{Consequence for rational \(5\)-tuples}

For odd \(n\), no positive constant vector lies in the image of \(P\). Indeed, if

\[
P(Y)=c\mathbf1,\qquad c>0,
\]

then there are signs \(s_i\in{\pm1}\) such that

\[
y_i-y_{i+1}=s_ic.
\]

Summing cyclically gives

\[
0=c\sum_{i=1}^n s_i,
\]

which is impossible when \(n\) is odd.

It follows that for odd \(n\), only constant starting vectors can ever reach zero exactly.

Therefore, for \(n=5\):

\[
\boxed{
X\in\mathbb Q_+^5,\ X\text{ nonconstant}
\quad\Longrightarrow\quad
\text{the orbit eventually enters the exact period-}15\text{ cycle}.
}
\]

\medskip\hrule\medskip

\subsection{8. Small and boundary cases}

\paragraph{Constant vectors}

If \(X=c\mathbf1\), then \(P(X)=0\).

\paragraph{\(n=3\)}

Here \(h=1\), so \(\mathcal S_3\) consists of the even-weight binary vectors:

\[
000,\quad110,\quad011,\quad101.
\]

The three nonzero vectors form the cycle

\[
110\longmapsto011\longmapsto101\longmapsto110.
\]

Every real \(3\)-orbit is asymptotic either to \(0\) or to a scalar multiple of this \(3\)-cycle.

\paragraph{\(n=4\)}

Here \(h=4\), so

\[
\mathcal S_4={0000}.
\]

Thus every real orbit converges to zero. This does not imply finite extinction: the algebraic exceptional family mentioned in the editorial note converges to zero without reaching it.

\paragraph{\(n=5\)}

There is one nonzero \(15\)-cycle up to scale and phase. Both nonperiodic possibilities occur:

\begin{itemize}
\item the eigenray of Section 5 approaches zero;
\item the construction of Section 6 approaches the nonzero \(15\)-cycle.
\end{itemize}

\paragraph{Necessity of the power-of-two hypothesis}

If \(n\) is not a power of \(2\), then

\[
|\mathcal S_n|=2^{n-h}>1,
\]

so nonzero periodic orbits exist. Hence ``every orbit converges to zero'' is equivalent to \(n\) being a power of \(2\).

\paragraph{Nonnegativity}

The restriction to \(\mathbb R_+^n\) is not essential: \(P\) sends every vector in \(\mathbb R^n\) into \(\mathbb R_+^n\) after one iteration.

\medskip\hrule\medskip

\subsection{9. Exact remaining obstruction}

The asymptotic shape is now determined for every \(n\), and completely explicit for \(n=5\). What is not supplied by Theorem 2.1 is a finite intrinsic criterion for an arbitrary real input \(X\) deciding whether

\[
c(X)=\lim_{q\to\infty}|P^q(X)|_\infty
\]

is zero or positive, nor a closed formula for \(c(X)\) and the eventual phase.

For \(n=5\), the exact obstruction can be formulated as follows.

Near the \(15\)-cycle, the \(15\)-step return map is piecewise linear. Each branch is a product of \(15\) signed difference matrices. A complete basin classification would require describing all admissible infinite paths through the resulting finite polyhedral branch graph and determining which associated matrix products contract the perturbation while preserving a positive limiting scale.

The construction in Section 6 gives one exact contracting return branch. A natural next lemma would be:

\textbf{Proposed next lemma.} Construct the finite directed graph of admissible polyhedral cones for the \(15\)-step return map at \(e_0\), and classify all strongly connected components whose matrix cocycle has a nontrivial stable direction.

That would turn the remaining \(n=5\) basin problem into a finite exact computation followed by a matrix-semigroup argument.

Thus the strongest justified status is:

\[
\boxed{
\begin{array}{l}
\text{Complete classification of all attractors for every }n;\\
\text{complete asymptotic classification for }n=5;\\
\text{exact rational-input algorithm;}\\
\text{explicit nonperiodic }n=5\text{ examples for both }c=0\text{ and }c>0;\\
\text{general real basin classification still remaining.}
\end{array}
}
\]

[1]: https://hrj.episciences.org/105/pdf "https://hrj.episciences.org/105/pdf"
[2]: https://www.fq.math.ca/Papers1/45-2/quartzerr02\_2007.pdf "https://www.fq.math.ca/Papers1/45-2/quartzerr02\_2007.pdf"
